\section{References}

Editions are identified as they were used. Where a work was read in a digital
transcription or text layer rather than in the printed edition itself, that is
stated. Loci given in the body are not repeated here.

\subsection{Scripture and versions}

\begin{itemize}[leftmargin=1.2em,itemsep=0.2em,topsep=0.3em]
\item \textit{The Holy Bible, Douay--Rheims}, Bishop Richard Challoner's revision
(1749--1752), in the electronic form registered in this repository. All English
scriptural quotations; Vulgate chapter and verse numbering throughout. Public
domain.
\item \textit{The Holie Bible faithfully translated into English, out of the
authentical Latin}, volume 1 (Douay, 1609). Read in the archive.org text layer of
the Princeton Theological Seminary copy. Public domain.
\item \textit{Biblia Sacra Vulgatae Editionis}, ed.\ Michael Hetzenauer, O.F.M.
Cap.\ (Ratisbon and Rome: Pustet, 1914). Genesis 15 read from a 200 dpi image of
numbered scan page 0012. Public domain.
\item Henry Barclay Swete, ed., \textit{The Old Testament in Greek according to
the Septuagint}, volume I, 4th ed.\ (Cambridge University Press, 1909; reprinted
1925). Genesis 15 text and apparatus read in the archive.org text layer. Public
domain.
\item Targum Onkelos and Targum Pseudo-Jonathan on Genesis 15. Aramaic read at
Sefaria; Pseudo-Jonathan also in J.~W.~Etheridge, \textit{The Targums of Onkelos
and Jonathan Ben Uzziel on the Pentateuch} (London, 1862). Public domain.
\end{itemize}

\subsection{Ancient witnesses}

\begin{itemize}[leftmargin=1.2em,itemsep=0.2em,topsep=0.3em]
\item Philo of Alexandria, \textit{Quis rerum divinarum heres sit}, in
\textit{The Works of Philo Judaeus}, trans.\ C.~D.~Yonge, volume II (London:
Bohn, 1854). Read in the archive.org text layer. Public domain.
\item Flavius Josephus, \textit{Antiquitates Judaicae} I.7.2, trans.\ William
Whiston, read in the University of Chicago (Penelope) transcription. Public
domain.
\item \textit{Babylonian Talmud}, tractates Shabbat and Pesachim. Vocalized
Aramaic and facing English of the William Davidson edition, at Sefaria. The
translation is licensed CC BY-NC and is not offered under this project's licence.
\item \textit{Bereshit Rabbah} 44:10 and 44:12. English of \textit{The Sefaria
Midrash Rabbah} (2022), licensed CC BY; Hebrew of the Torat Emet text.
\item Augustine, \textit{De civitate Dei} XVI.23--24. Latin in Emanuel Hoffmann,
ed., \textit{CSEL} 40 (Vienna, 1899--1900); English in Marcus Dods, trans.,
\textit{The City of God} (Edinburgh, 1871). Both registered in this repository;
both public domain.
\end{itemize}

\subsection{Jewish commentary}

\begin{itemize}[leftmargin=1.2em,itemsep=0.2em,topsep=0.3em]
\item Rashi on Genesis. \textit{Pentateuch with Rashi's Commentary}, trans.\
M.~Rosenbaum and A.~M.~Silbermann (1929--1934). Public domain.
\item Abraham Ibn Ezra on Genesis. Hebrew of the Piotrkow edition (1907--1911),
public domain; English of H.~Norman Strickman and Arthur M.~Silver (Menorah,
1988--2004), licensed CC BY-NC and not offered under this project's licence.
\item Moses Nahmanides (Ramban) on Genesis. \textit{Commentary on the Torah},
trans.\ Charles B.~Chavel (New York: Shilo, 1971--1976), licensed CC BY.
\item David Kimhi (Radak) on Genesis, in Eliyahu Munk, \textit{HaChut
Hameshulash}, licensed CC BY.
\item Hezekiah ben Manoah (Chizkuni) on Genesis. Hebrew public domain; English of
Eliyahu Munk, licensed CC BY.
\item Bahya ben Asher, \textit{Midrash Rabbeinu Bachya} (Warsaw, 1878). Hebrew
public domain; English of Eliyahu Munk, licensed CC BY.
\item Jacob ben Asher, \textit{Tur HaArokh} on Genesis (Hanover, 1838). Hebrew
public domain.
\item Meir Leibush ben Yehiel Michel (Malbim) on Genesis, Hebrew Wikisource text.
Public domain.
\item Naftali Zvi Yehuda Berlin (Netziv), \textit{Haamek Davar}, in \textit{Sefer
Torat Elohim} (Vilna, 1879). Public domain.
\item Isaac Abarbanel, \textit{Perush al ha-Torah}, on Genesis (Warsaw, 1862).
Read in two independent digital transcriptions, at Sefaria and at Hebrew
Wikisource, and collated between them for the passage quoted. Public domain.
\item Moses Maimonides, \textit{The Guide for the Perplexed}, trans.\
M.~Friedl\"ander, 2nd ed.\ (London, 1903). Public domain.
\end{itemize}

\subsection{Catholic commentary}

\begin{itemize}[leftmargin=1.2em,itemsep=0.2em,topsep=0.3em]
\item Cornelius a Lapide, S.J., \textit{Commentaria in Pentateuchum Mosis}, on
Genesis. Read in the archive.org text layer of a 1700 printing. Public domain.
\item George Leo Haydock, notes to \textit{The Douay--Rheims Edition of the Holy
Catholic Bible with a Comprehensive Catholic Commentary}. Read in the archive.org
text layer. Public domain.
\end{itemize}

\subsection{On the digital transcriptions used}

Several of the Jewish sources above were read through Sefaria's digital library,
which supplies both the Hebrew or Aramaic base text and, where one exists, a
facing translation. Sefaria is a delivery route, not an edition: in every case
the edition it serves is named in the list above and in the witness register, and
the rights status recorded is that of the named edition and translation, not of
the platform. Where the platform's own base text has no named printed exemplar,
that is stated as a limitation in the preceding appendix rather than concealed
behind the platform's name.
